\documentclass[../../main.tex]{subfiles}
\begin{document}

\section{坐标中的计算}\label{section:ch03-坐标中的计算}

到目前为止，我们对一个流形切空间的处理也许看起来抽象得令人绝望。为把它落到实处，我们将说明如何在局部坐标中用切向量和微分作计算。

首先，设 $M$ 是一个光滑（无边）流形，$(U,\varphi)$ 是 $M$ 上的一个光滑坐标图。则特别地，$\varphi$ 是从 $U$ 到开子集 $\widetilde{U}\subset\mathbb{R}^{n}$ 的微分同胚。结合命题~\cref{proposition:ch03-开子流形切空间} 和命题~\cref{proposition:ch03-微分性质}(d)，我们看到 $\mathrm{d}\varphi_{p}:T_{p}M\to T_{\varphi(p)}\mathbb{R}^{n}$ 是一个同构。

由推论~\cref{corollary:ch03-偏导基}，导子 $\partial/\partial x^{1}|_{\varphi(p)},\dots,\partial/\partial x^{n}|_{\varphi(p)}$ 构成 $T_{\varphi(p)}\mathbb{R}^{n}$ 的一组基。因此，这些向量在同构 $\mathrm{d}\varphi_{p}$ 下的原像构成 $T_{p}M$ 的一组基（图~\cref{figure:ch03-3.5}）。按照我们在可能时把坐标映射视作等同的惯例，我们用记号 $\partial/\partial x^{i}|_{p}$ 表示这些向量，其特征由下列表达式之一给出：
\begin{align}\label{align:ch03-坐标向量定义}
\left.\frac{\partial}{\partial x^{i}}\right|_{p}
=(\mathrm{d}\varphi_{p})^{-1}\left(\left.\frac{\partial}{\partial x^{i}}\right|_{\varphi(p)}\right)
=\mathrm{d}(\varphi^{-1})_{\varphi(p)}\left(\left.\frac{\partial}{\partial x^{i}}\right|_{\varphi(p)}\right).
\end{align}

\begin{figure}[ht]
\centering
\includegraphics[width=0.8\textwidth]{fig-03-5.png}
\caption{坐标中的切向量}\label{figure:ch03-3.5}
\end{figure}

展开定义，我们看到 $\partial/\partial x^{i}|_{p}$ 对函数 $f\in C^{\infty}(U)$ 的作用为
\begin{align*}
\left.\frac{\partial}{\partial x^{i}}\right|_{p}f
=\left.\frac{\partial}{\partial x^{i}}\right|_{\varphi(p)}(f\circ\varphi^{-1})
=\frac{\partial\widetilde{f}}{\partial x^{i}}(\widetilde{p}),
\end{align*}
其中 $\widetilde{f}=f\circ\varphi^{-1}$ 是 $f$ 的坐标表示，$\widetilde{p}=(p^{1},\dots,p^{n})=\varphi(p)$ 是 $p$ 的坐标表示。换言之，$\partial/\partial x^{i}|_{p}$ 正是那个在 $p$ 的坐标表示处、取 $f$ 的坐标表示对 $x^{i}$ 的第 $i$ 个偏导数的导子。向量 $\partial/\partial x^{i}|_{p}$ 称为与给定坐标系相关的 $p$ 处的坐标向量。在 $\mathbb{R}^{n}$ 上标准坐标的特殊情形，向量 $\partial/\partial x^{i}|_{p}$ 字面上就是偏导数算子。

当 $M$ 是一个光滑带边流形而 $p$ 是一个内点时，上面的讨论逐字适用。对 $p\in\partial M$，唯一需要做的改变是把 $\mathbb{R}^{n}$ 换成 $\mathbb{H}^{n}$，并理解记号 $\partial/\partial x^{i}|_{\varphi(p)}$ 可以互换地表示 $T_{\varphi(p)}\mathbb{R}^{n}$ 的一个元素或 $T_{\varphi(p)}\mathbb{H}^{n}$ 的一个元素，这符合我们把同构 $\mathrm{d}\iota_{\varphi(p)}:T_{\varphi(p)}\mathbb{H}^{n}\to T_{\varphi(p)}\mathbb{R}^{n}$ 视作等同的约定。此时第 $n$ 个坐标向量 $\partial/\partial x^{n}|_{p}$ 应解释为单侧导数。

下面的命题总结了到目前为止的讨论。

\begin{proposition}\label{proposition:ch03-坐标向量基}
设 $M$ 是一个 $n$ 维光滑带边或无边流形，$p\in M$。则 $T_{p}M$ 是一个 $n$ 维向量空间，且对任意包含 $p$ 的光滑坐标图 $\bigl(U,(x^{i})\bigr)$，坐标向量 $\partial/\partial x^{1}|_{p},\dots,\partial/\partial x^{n}|_{p}$ 构成 $T_{p}M$ 的一组基。
\end{proposition}

于是，一个切向量 $v\in T_{p}M$ 可以唯一地写成线性组合
\begin{align*}
v=v^{i}\left.\frac{\partial}{\partial x^{i}}\right|_{p},
\end{align*}
其中我们照例使用求和约定，且如第~52 页所解释的，分母中的上指标被看作下指标。有序基 $(\partial/\partial x^{i}|_{p})$ 称为 $T_{p}M$ 的一组坐标基，数 $(v^{1},\dots,v^{n})$ 称为 $v$ 关于该坐标基的分量。若 $v$ 已知，其分量可以容易地从它在坐标函数上的作用算出。对每个 $j$，$v$ 的分量由 $v^{j}=v(x^{j})$ 给出（这里我们把 $x^{j}$ 看作 $U$ 上的光滑实值函数），因为
\begin{align*}
v(x^{j})=\left(v^{i}\left.\frac{\partial}{\partial x^{i}}\right|_{p}\right)(x^{j})
=v^{i}\frac{\partial x^{j}}{\partial x^{i}}(p)=v^{j}.
\end{align*}

\subsection{坐标中的微分}\label{subsection:ch03-坐标中的微分}

接着我们考察微分在坐标中是什么样子。我们从考虑一个光滑映射 $F:U\to V$ 的特殊情形开始，其中 $U\subset\mathbb{R}^{n}$、$V\subset\mathbb{R}^{m}$ 是欧几里得空间的开子集。对任意 $p\in U$，我们将用标准坐标基确定 $\mathrm{d}F_{p}:T_{p}\mathbb{R}^{n}\to T_{F(p)}\mathbb{R}^{m}$ 的矩阵。用 $(x^{1},\dots,x^{n})$ 记定义域中的坐标，$(y^{1},\dots,y^{m})$ 记陪域中的坐标，我们用链式法则计算 $\mathrm{d}F_{p}$ 在一个典型基向量上的作用如下：
\begin{align*}
\mathrm{d}F_{p}\left(\left.\frac{\partial}{\partial x^{i}}\right|_{p}\right)f
=\left.\frac{\partial}{\partial x^{i}}\right|_{p}(f\circ F)
=\frac{\partial f}{\partial y^{j}}\bigl(F(p)\bigr)\frac{\partial F^{j}}{\partial x^{i}}(p)
=\left(\frac{\partial F^{j}}{\partial x^{i}}(p)\left.\frac{\partial}{\partial y^{j}}\right|_{F(p)}\right)f.
\end{align*}
因此
\begin{align}\label{align:ch03-欧氏微分基}
\mathrm{d}F_{p}\left(\left.\frac{\partial}{\partial x^{i}}\right|_{p}\right)
=\frac{\partial F^{j}}{\partial x^{i}}(p)\left.\frac{\partial}{\partial y^{j}}\right|_{F(p)}.
\end{align}
换言之，$\mathrm{d}F_{p}$ 在坐标基下的矩阵是
\begin{align*}
\begin{pmatrix}
\frac{\partial F^{1}}{\partial x^{1}}(p) & \cdots & \frac{\partial F^{1}}{\partial x^{n}}(p)\\
\vdots & \ddots & \vdots\\
\frac{\partial F^{m}}{\partial x^{1}}(p) & \cdots & \frac{\partial F^{m}}{\partial x^{n}}(p)
\end{pmatrix}.
\end{align*}
（回顾矩阵的列是基向量像的分量。）这个矩阵正是 $F$ 在 $p$ 处的 Jacobian 矩阵，它是全导数 $DF(p):\mathbb{R}^{n}\to\mathbb{R}^{m}$ 的矩阵表示。因此，在这个情形，在我们把欧几里得空间与其切空间通常等同下，$\mathrm{d}F_{p}:T_{p}\mathbb{R}^{n}\to T_{F(p)}\mathbb{R}^{m}$ 对应于全导数 $DF(p):\mathbb{R}^{n}\to\mathbb{R}^{m}$。若 $U$ 是 $\mathbb{H}^{n}$ 的开子集、$V$ 是 $\mathbb{H}^{m}$ 的开子集，同样的计算适用。

\begin{figure}[ht]
\centering
\includegraphics[width=0.8\textwidth]{fig-03-6.png}
\caption{坐标中的微分}\label{figure:ch03-3.6}
\end{figure}

现在考虑带边或无边光滑流形之间光滑映射 $F:M\to N$ 这一更一般的情形。选取 $M$ 包含 $p$ 的光滑坐标图 $(U,\varphi)$ 和 $N$ 包含 $F(p)$ 的光滑坐标图 $(V,\psi)$，我们得到坐标表示 $\widetilde{F}=\psi\circ F\circ\varphi^{-1}:\varphi(U\cap F^{-1}(V))\to\psi(V)$（图~\cref{figure:ch03-3.6}）。令 $\widetilde{p}=\varphi(p)$ 记 $p$ 的坐标表示。由上面的计算，$\mathrm{d}\widetilde{F}_{\widetilde{p}}$ 关于标准坐标基由 $\widetilde{F}$ 在 $\widetilde{p}$ 处的 Jacobian 矩阵表示。利用 $F\circ\varphi^{-1}=\psi^{-1}\circ\widetilde{F}$ 这一事实，我们计算
\begin{align}\label{align:ch03-流形微分基}
\mathrm{d}F_{p}\left(\left.\frac{\partial}{\partial x^{i}}\right|_{p}\right)
=\mathrm{d}F_{p}\left(\mathrm{d}(\varphi^{-1})_{\widetilde{p}}\left(\left.\frac{\partial}{\partial x^{i}}\right|_{\widetilde{p}}\right)\right)
=\mathrm{d}(\psi^{-1})_{\widetilde{F}(\widetilde{p})}\left(\mathrm{d}\widetilde{F}_{\widetilde{p}}\left(\left.\frac{\partial}{\partial x^{i}}\right|_{\widetilde{p}}\right)\right)
=\frac{\partial\widetilde{F}^{j}}{\partial x^{i}}(\widetilde{p})\left.\frac{\partial}{\partial y^{j}}\right|_{F(p)}.
\end{align}
因此，$\mathrm{d}F_{p}$ 在坐标基下由 $F$ 的坐标代表的 Jacobian 矩阵表示。事实上，微分的定义正是为了给 Jacobian 矩阵一个与坐标无关的意义而量身定做的。

在微分几何文献中，微分有时称为切映射、全导数，或简称 $F$ 的导数。由于它把切向量"前推"，从定义域流形带到陪域，它也称为（逐点）前推。不同作者用诸如 $F'(p)$、$DF$、$DF(p)$、$F_{*}$、$TF$、$T_{p}F$ 等记号表示它。我们将坚持用记号 $\mathrm{d}F_{p}$ 表示流形之间光滑映射的微分，而把 $DF(p)$ 留给有限维向量空间之间映射的全导数，在欧几里得空间情形我们把它等同于 $F$ 的 Jacobian 矩阵。

\subsection{坐标变换}\label{subsection:ch03-坐标变换}

设 $(U,\varphi)$ 和 $(V,\psi)$ 是 $M$ 上的两个光滑坐标图，$p\in U\cap V$。我们用 $(x^{i})$ 记 $\varphi$ 的坐标函数，用 $(\widetilde{x}^{i})$ 记 $\psi$ 的坐标函数。$p$ 处的任意切向量可以关于基 $(\partial/\partial x^{i}|_{p})$ 或 $(\partial/\partial\widetilde{x}^{i}|_{p})$ 表示。这两种表示如何关联？

在这种情形，习惯上用如下简写记号写过渡映射 $\psi\circ\varphi^{-1}:\varphi(U\cap V)\to\psi(U\cap V)$：
\begin{align*}
\psi\circ\varphi^{-1}(x)=\bigl(\widetilde{x}^{1}(x),\dots,\widetilde{x}^{n}(x)\bigr).
\end{align*}
这里我们纵容一种典型的记号滥用：在表达式 $\widetilde{x}^{i}(x)$ 中，我们把 $\widetilde{x}^{i}$ 看作一个坐标函数（其定义域是 $M$ 的一个开子集，等同于 $\mathbb{R}^{n}$ 或 $\mathbb{H}^{n}$ 的一个开子集）；但我们把 $x$ 看作代表一个点（此处即在 $\varphi(U\cap V)$ 中）。由~\eqref{align:ch03-欧氏微分基}，微分 $\mathrm{d}(\psi\circ\varphi^{-1})_{\varphi(p)}$ 可写成
\begin{align*}
\mathrm{d}(\psi\circ\varphi^{-1})_{\varphi(p)}\left(\left.\frac{\partial}{\partial x^{i}}\right|_{\varphi(p)}\right)
=\frac{\partial\widetilde{x}^{j}}{\partial x^{i}}\bigl(\varphi(p)\bigr)\left.\frac{\partial}{\partial\widetilde{x}^{j}}\right|_{\psi(p)}.
\end{align*}
利用坐标向量的定义，我们得到
\begin{align}\label{align:ch03-坐标向量变换}
\left.\frac{\partial}{\partial x^{i}}\right|_{p}
=\mathrm{d}(\psi^{-1})_{\psi(p)}\circ\mathrm{d}(\psi\circ\varphi^{-1})_{\varphi(p)}\left(\left.\frac{\partial}{\partial x^{i}}\right|_{\varphi(p)}\right)
=\frac{\partial\widetilde{x}^{j}}{\partial x^{i}}(\widetilde{p})\left.\frac{\partial}{\partial\widetilde{x}^{j}}\right|_{p},
\end{align}
其中我们再次写 $\widetilde{p}=\varphi(p)$。（这个公式容易记忆，因为它看起来与 $\mathbb{R}^{n}$ 中偏导数的链式法则完全相同。）把它应用于向量 $v=v^{i}\partial/\partial x^{i}|_{p}=\widetilde{v}^{j}\partial/\partial\widetilde{x}^{j}|_{p}$ 的分量，我们发现 $v$ 的分量按规则
\begin{align}\label{align:ch03-分量变换}
\widetilde{v}^{j}=\frac{\partial\widetilde{x}^{j}}{\partial x^{i}}(\widetilde{p})v^{i}
\end{align}
变换。

\begin{figure}[ht]
\centering
\includegraphics[width=0.8\textwidth]{fig-03-7.png}
\caption{坐标变换}\label{figure:ch03-3.7}
\end{figure}

\begin{example}[极坐标]\label{example:ch03-极坐标}
平面中适当开子集上极坐标与标准坐标之间的过渡映射由 $(x,y)=(r\cos\theta,r\sin\theta)$ 给出。令 $p$ 为 $\mathbb{R}^{2}$ 中极坐标表示为 $(r,\theta)=(2,\pi/2)$ 的点，令 $v\in T_{p}\mathbb{R}^{2}$ 为极坐标表示为
\begin{align*}
v=3\left.\frac{\partial}{\partial r}\right|_{p}-\left.\frac{\partial}{\partial\theta}\right|_{p}
\end{align*}
的切向量。把~\eqref{align:ch03-坐标向量变换} 应用于坐标向量，我们求得
\begin{align*}
\left.\frac{\partial}{\partial r}\right|_{p}=\cos\frac{\pi}{2}\left.\frac{\partial}{\partial x}\right|_{p}+\sin\frac{\pi}{2}\left.\frac{\partial}{\partial y}\right|_{p}=\left.\frac{\partial}{\partial y}\right|_{p},
\end{align*}
\begin{align*}
\left.\frac{\partial}{\partial\theta}\right|_{p}=-2\sin\frac{\pi}{2}\left.\frac{\partial}{\partial x}\right|_{p}+2\cos\frac{\pi}{2}\left.\frac{\partial}{\partial y}\right|_{p}=-2\left.\frac{\partial}{\partial x}\right|_{p},
\end{align*}
从而 $v$ 在标准坐标中有如下坐标表示：
\begin{align*}
v=3\left.\frac{\partial}{\partial y}\right|_{p}+2\left.\frac{\partial}{\partial x}\right|_{p}.
\end{align*}
\end{example}

需要牢记的一个重要事实是，每个坐标向量 $\partial/\partial x^{i}|_{p}$ 依赖于整个坐标系，而不仅依赖于单个坐标函数 $x^{i}$。几何上，这反映了如下事实：$\partial/\partial x^{i}|_{p}$ 是在所有其他坐标保持不变时对 $x^{i}$ 求导所得的导子。若除 $x^{i}$ 外的坐标函数被改变，则这个坐标导数的方向可能改变。下一个练习说明这如何发生。

\begin{exercise}\label{exercise:ch03-3.17}
令 $(x,y)$ 记 $\mathbb{R}^{2}$ 上的标准坐标。验证 $(\widetilde{x},\widetilde{y})$ 是 $\mathbb{R}^{2}$ 上的整体光滑坐标，其中
\begin{align*}
\widetilde{x}=x,\qquad \widetilde{y}=y+x^{3}.
\end{align*}
令 $p$ 为点 $(1,0)\in\mathbb{R}^{2}$（在标准坐标中），并证明
\begin{align*}
\left.\frac{\partial}{\partial x}\right|_{p}\neq\left.\frac{\partial}{\partial\widetilde{x}}\right|_{p},
\end{align*}
尽管坐标函数 $x$ 与 $\widetilde{x}$ 恒等。
\end{exercise}

\end{document}